% The running example's 17 x 19 output against the 16 x 16 tiles that cover it.
\begin{figure}[htbp]
\centering
\begin{tikzpicture}[x=1mm, y=1mm,
    lbl/.style={tag, font=\scriptsize, align=center},
    num/.style={circle, draw=ink!70, fill=white, inner sep=0.8pt, font=\scriptsize, text=ink},
    key/.style={font=\scriptsize, text=ink, anchor=west}]

  \def\u{1.9}   % one element
  \def\g{60.8}  % 32 elements: the two tile rows and columns together

  % The part of the tile grid the matrix actually reaches.
  \fill[ours!30] (0,0) rectangle ({19*\u},{-17*\u});

  % The whole 32 x 32 tile grid, and the tile boundaries inside it.
  \draw[ink!35, line width=0.4pt, dashed] (0,0) rectangle ({\g},{-\g});
  \draw[ink!35, line width=0.4pt, dashed] ({16*\u},0) -- ({16*\u},{-\g});
  \draw[ink!35, line width=0.4pt, dashed] (0,{-16*\u}) -- ({\g},{-16*\u});

  % The live result on top of it.
  \draw[ink!75, line width=0.7pt] (0,0) rectangle ({19*\u},{-17*\u});
  \draw[ink!75, line width=0.7pt] ({16*\u},0) -- ({16*\u},{-17*\u});
  \draw[ink!75, line width=0.7pt] (0,{-16*\u}) -- ({19*\u},{-16*\u});

  \node[num] at ({8*\u},{-8*\u})  {1};
  \node[num] at ({24*\u},{-8*\u}) {2};
  \node[num] at ({8*\u},{-24*\u}) {3};
  \node[num] at ({24*\u},{-24*\u}){4};

  \draw[decorate, decoration={brace, amplitude=4pt}, line width=0.5pt, color=ink!70]
    (0,1.5) -- ({19*\u},1.5) node[midway, above=4pt, lbl] {19 columns};
  \draw[decorate, decoration={brace, amplitude=4pt, mirror}, line width=0.5pt, color=ink!70]
    (-1.5,0) -- (-1.5,{-17*\u}) node[midway, left=4pt, lbl] {17\\rows};

  % Key, clear of the drawing.
  \node[key] at ({\g+9},-4)  {\textbf{1}\ \ 16 × 16, the only full tile\ \ R24--R31};
  \node[key] at ({\g+9},-11) {\textbf{2}\ \ 16 × 3\ \ R32--R39};
  \node[key] at ({\g+9},-18) {\textbf{3}\ \ 1 × 16\ \ R40--R47};
  \node[key] at ({\g+9},-25) {\textbf{4}\ \ 1 × 3\ \ R48--R55};
  \node[key, text width=52mm, align=left, text=ink!65] at ({\g+9},-36)
    {Shaded: the result. Dashed: the rest of each 16 × 16 tile, which the 20 mask instructions
     switch off.};
\end{tikzpicture}
\caption{The running example's output against the tiles that cover it. A 17 × 19 result does not divide into
16 × 16, so it takes two tile rows by two tile columns, three of them partial: one \op{5107} and one
eight-register accumulator group each.}
\label{fig:example-tiles}
\end{figure}
